\section{What an isoperimetric theorem means in intuitionistic PRA}
\label{sec:foundations}

The first-order language of \(\ipra\) has natural-number variables, symbols
for primitive recursive functions, and intuitionistic logical operations.
It does not have a sort containing all subsets of \(\mathbb R^n\), all
continuous functions, or all real numbers. Accordingly, there are two
precise kinds of result in this report.

\begin{enumerate}
\item A \emph{uniform finite theorem} has arbitrary natural-number inputs
encoding rational numbers, finite lists, graphs, or boxes. Its hypotheses
and conclusion are decidable comparisons, or explicitly displayed
first-order formulas built from them.
\item An \emph{effective geometric theorem schema} fixes primitive
recursive names and moduli, possibly with additional numerical parameters.
The theorem then concerns all those parameters and all requested
precisions. A name is never evaluated by a supposed universal total
evaluator for arbitrary primitive recursive programs.
\end{enumerate}

This distinction is essential. There is no primitive recursive universal
evaluator for all primitive recursive programs: a diagonal program would
contradict its universality. Fixing a family of programs, or accepting
bounded computation traces, is legitimate; silently quantifying over
arbitrary functions and evaluating them is not.

\subsection{Coding and order}

Encode an integer by its sign and magnitude and a rational by an integer
numerator and a positive denominator. Comparisons use cross products;
normalization by the Euclidean algorithm is optional. Finite lists have
a length and bounded access functions. Products, sums, bit operations,
bounded minimization, and finite sorting are primitive recursive.
The displayed rational inequalities can therefore be cleared of positive
denominators and read as integer inequalities.

A named real \(x\) is a fixed primitive recursive rational sequence
\(x[k]\) with a proved Cauchy estimate, for example
\[
 |x[j]-x[k]|\leq 2^{-j}+2^{-k}\quad(j,k\in\mathbb N).
\]
The intended value satisfies \(|x-x[k]|\leq2^{-k}\). This latter notation
is explanatory shorthand: every argument can instead use the preceding
rational comparisons. Define
\[
 x\leq y\quad:\Longleftrightarrow\quad
 (\forall k)\ x[k]\leq y[k]+2^{1-k},
\]
after adopting synchronized error bounds. Equality means both orders,
or an equivalent zero-distance condition. A positive real used as a
denominator is accompanied by a rational lower bound \(0<c\leq x\).
There is no case distinction between \(x=0\) and \(x>0\) for an arbitrary
named real. All case distinctions below concern integers, rational numbers,
or supplied positive-apartness evidence.

If an argument proves \(x\leq y+2^{-k}\) for every \(k\), then it proves
\(x\leq y\) by this definition of order. It does not invoke a compactness
or least-upper-bound principle. The zero conclusions used in the equality
proofs have the same form: a nonnegative term bounded by arbitrarily small
dyadics is equal to zero in the representation.

\subsection{Elementary numerical operations with bounds}

For rational \(0\leq x\leq M\) and \(m\geq1\), bisection on
\([0,M+1]\), comparing rational \(m\)-th powers, produces lower and
upper bounds on \(x^{1/m}\). Bisection a prescribed number of times is
primitive recursion. On nonnegative inputs,
\[
 |x^{1/m}-y^{1/m}|\leq |x-y|^{1/m};
\]
indeed \((a+b)^m\geq a^m+b^m\) for \(a,b\geq0\). For a named
nonnegative input, first obtain a rational enclosure of width at most
\(2^{-m(k+3)}\), truncating its lower endpoint at zero. Apply the
rational bisection routine separately to its two endpoints, with endpoint
root errors at most \(2^{-k-3}\). The displayed inequality bounds the
gap between the exact endpoint roots by \(2^{-k-3}\), so the resulting
rational interval gives the required root precision. Every branching
comparison is rational; one never compares a rational midpoint with an
arbitrary named real by a supposedly decidable order test.
For \(0\leq x,y\leq B\),
\[
 |x^m-y^m|\leq mB^{m-1}|x-y|
\]
follows by factoring the difference. These two estimates give explicit
precision budgets for every fractional power in the report.

The reused planar appendix constructs \(\pi\), sine and cosine from
dyadic regular polygons and positive quadratic radicals, with quantitative
estimates. Alternatively, rational Taylor sums on a prescribed bounded
interval, with the factorial tail bounded in advance, give the same
primitive recursive names. Only those named constants and operations are
needed here. No unrestricted theory of real analysis is assumed.

\subsection{Why the finite inductions stay inside \(\ipra\)}

\(\ipra\) permits induction for quantifier-free formulas. A bounded
universal statement \(\forall i<B\,R(i)\), with primitive recursive
decidable \(R\), is equivalent to the equation saying that the product
of its Boolean characteristic values is one. Bounded existential
statements are handled by a bounded sum or scan. Their elementary
equivalences are themselves proved by quantifier-free induction.

This observation is used with actual bounds, not as a substitute for them.
The graph proof supplies bounds on masks and on the number of compression
steps. The box proof supplies a decreasing fuel and a primitive recursive
bound on every intermediate rational. The planar appendix bounds the
number of subdivisions, Fourier terms, and radical approximations. A
finite proof tree can consequently be traversed in postorder; induction
asserts that every previously processed node has the claimed rational
inequality. This is a bounded Boolean invariant.

There is a further logical restriction. Proving a comparison about a
\emph{particular} rational input by computation is not a proof of the
uniform theorem. The human arguments below prove preservation of these
invariants for arbitrary numerical inputs. The executable examples are
regression checks for the formulas and algorithms; they are not the
justification for their universality.

\subsection{How to read the formalization status}

The mathematical arguments here are non-formalized proofs of the stated
finite theorems and effective schemas, with their arithmetic compilation
explained. The accompanying Lean material is a development blueprint.
No claim is made that the present archive contains completed Lean terms
of all the relevant internal provability types. In particular, a Mathlib
proof using \(\mathbb R\), choice, or general measure theory would not
by itself establish a theorem in \(\ipra\). The repository audit in
Section~\ref{sec:repository-audit} records the actual bridge that must
be used after an internal derivation has been constructed.

